\documentclass[11pt]{article}
\usepackage{amsmath,amsthm,amssymb}
\usepackage[margin=1in]{geometry}
\usepackage{hyperref}

\newtheorem{theorem}{Theorem}
\newtheorem{lemma}{Lemma}
\newtheorem{proposition}{Proposition}
\newtheorem*{remark}{Remark}

\title{Structural proof of the column-1 SC discriminator for $\Omega$ on $\lambda=(3,3,1,1,1)$}
\author{Clio}
\date{2026-05-19}

\newcommand{\Vp}{V_{+}}
\newcommand{\Vm}{V_{-}}
\newcommand{\Vpalpha}{V_{+}^{\alpha}}
\newcommand{\Vpbeta}{V_{+}^{\beta}}

\begin{document}
\maketitle

\section{Setup}

Let $\lambda = (3,3,1,1,1)$, $n = 9$. Working in the Hecke algebra $H_n(q)$
with Hoefsmit's semi-normal basis $\{v_T : T \in \mathrm{SYT}(\lambda)\}$, $N = 120$.
The generators are $T_i$ ($i=1,\ldots,8$); set $S_i := T_i + 1$.
Define
\[
R'_k := S_{k-1} S_{k-2} \cdots S_2 S_1 \qquad (k \ge 2),
\qquad
\Omega := R'_3 R'_4 R'_5 R'_6 R'_7 R'_8 R'_9.
\]
For a tableau $T$, write $T(r,c)$ for the entry in row $r$, column $c$ (1-indexed).

For each SYT $T$, $T_1$ has the eigenvalues $\{q,-1\}$ on $v_T$:
$T_1 v_T = q v_T$ when $T(1,1)=1$, $T(1,2)=2$ (\emph{SR-at-1}), and $T_1 v_T = -v_T$
when $T(1,1)=1$, $T(2,1)=2$ (\emph{SC-at-1}). Every SYT of $\lambda$ is one or the
other. Hence $V = \Vp \oplus \Vm$ where
\[
\Vp := \mathrm{span}\{v_T : T \text{ SR-at-1}\}, \qquad
\Vm := \mathrm{span}\{v_T : T \text{ SC-at-1}\}.
\]
$S_1 = T_1 + 1$ satisfies $S_1 v = (q+1) v$ on $\Vp$ and $S_1 v = 0$ on $\Vm$.
So $\Vm = \ker S_1$.

\paragraph{Three structural conditions on SR-at-1 SYTs.}
\begin{align*}
\text{(}\alpha\text{)} \quad & T(2,1) = 3 \text{ and } T(3,1) = 4 \\
\text{(}\beta\text{)} \quad & T(3,1) = 5 \text{ and } T(4,1) = 6 \\
\text{(}\gamma\text{)} \quad & T(4,1) = 7 \text{ and } T(5,1) = 8
\end{align*}
Each condition is a \emph{column-1 SC pair} at indices 3, 5, 7 respectively.
$\alpha$-SYTs form a 20-dimensional subspace
$\Vpalpha = \Vp \cap \ker S_3$ inside $\Vp$ (since in our shape, the only way 3,4 can
be column-adjacent in an SR-at-1 SYT is at $(2,1),(3,1)$). Similarly $\beta$-SYTs span
$\Vpbeta$.

\section{Theorem}

\begin{theorem} \label{thm:main}
Let $T \in \mathrm{SYT}(\lambda)$ be SR-at-1. Then $\Omega v_T = 0$ if and only if
$T$ satisfies at least one of {\normalfont(α), (β), (γ)}.
\end{theorem}

The converse direction is a finite check on the $50$ SR-at-1 SYTs (no $T$ outside
$\alpha \cup \beta \cup \gamma$ has $\Omega v_T = 0$). The structural content
is the forward direction: three implications $\alpha \Rightarrow \Omega v_T = 0$,
$\beta \Rightarrow \Omega v_T = 0$, $\gamma \Rightarrow \Omega v_T = 0$,
which we prove as Lemmas \ref{lem:Lalpha}, \ref{lem:Lbeta}, \ref{lem:Lgamma}.

\section{Key reduction}

Every $R'_k$ has $S_1$ as its rightmost factor, so $R'_k w = 0$ whenever
$S_1 w = 0$, i.e., $R'_k(\Vm) = 0$. In particular, to show $\Omega v_T = 0$
it suffices to show that at some stage $k \in \{3,\ldots,8\}$ we have
$\Pi_{k+1}(T) := R'_{k+1} R'_{k+2} \cdots R'_9 v_T \in \Vm$.

\subsection*{The $\Vp/\Vm$ decomposition and operator action}

\begin{lemma} \label{lem:preserves}
For $j \ge 3$, $S_j$ preserves $\Vp$ and $\Vm$ separately.
\end{lemma}
\begin{proof}
$S_j v_{T'} = a v_{T'} + b v_{s_j T'}$ in Hoefsmit basis (or a diagonal action),
where $s_j T'$ swaps the values $j, j+1$ in $T'$. For $j \ge 3$, this does not
move the entries $1$ or $2$, so $T'$ and $s_j T'$ are in the same $\Vp/\Vm$ class.
\end{proof}

\begin{lemma} \label{lem:S1S2S1}
If $T$ is SR-at-1 with $T(2,1) = 3$, then $S_1 S_2 S_1 v_T = -(q+1) v_T$.
\end{lemma}
\begin{proof}
$S_1 v_T = (q+1) v_T$ by SR-at-1. With $2$ at $(1,2)$ (content $1$) and $3$ at
$(2,1)$ (content $-1$), the axial distance is $d_2 = -2$, so
$S_2 v_T = a v_T + b v_{s_2 T}$ with $a = (q-1) q^{-2}/(q^{-2}-1) = -1/(q+1)$
(and some $b$). In $s_2 T$: 1 at $(1,1)$, 2 at $(2,1)$ — \emph{SC-at-1}, so
$v_{s_2 T} \in \Vm$ and $S_1 v_{s_2 T} = 0$.
Therefore
$
S_1 S_2 S_1 v_T = (q+1) (a (q+1) v_T + b \cdot 0) = a(q+1)^2 v_T = -(q+1) v_T.
$
\end{proof}

\section{Proof of L\textsubscript{α}}

\begin{lemma}[$L_\alpha$] \label{lem:Lalpha}
If $T$ is an $\alpha$-SYT, then $R'_k v_T \in \Vm$ for every $k \ge 4$.
In particular $\Omega v_T = 0$.
\end{lemma}

\begin{proof}
By Lemma~\ref{lem:S1S2S1}, $S_1 S_2 S_1 v_T = -(q+1) v_T$. Since $S_1$ commutes
with $S_j$ for $j \ge 3$, we may pull $S_1$ rightward in $S_1 R'_k =
S_1 \cdot S_{k-1} S_{k-2} \cdots S_1$:
\[
S_1 R'_k = (S_{k-1} \cdots S_3) \cdot S_1 S_2 S_1.
\]
Applied to $v_T$:
\[
S_1 R'_k v_T \;=\; (S_{k-1} \cdots S_3) \cdot \bigl(-(q+1) v_T\bigr)
\;=\; -(q+1) (S_{k-1} \cdots S_4) \cdot S_3 v_T.
\]
By the $\alpha$-condition $T(2,1)=3, T(3,1)=4$, entries $3,4$ are column-adjacent
in $T$, so $T_3 v_T = -v_T$ and $S_3 v_T = 0$. Therefore $S_1 R'_k v_T = 0$,
i.e.\ $R'_k v_T \in \Vm$. Since $R'_j \Vm = 0$ for every $j$, applying any further
$R'_j$ for $j < k$ kills it, giving $\Omega v_T = 0$.
\end{proof}

\section{Proof of L\textsubscript{β}}

The proof works by showing that the $\Vp$-part of $R'_9 v_T$ lands on $\alpha$-SYTs,
then applying $L_\alpha$ per-tableau.

\begin{lemma}[Support lemma SL\textsubscript{β}] \label{lem:SLbeta}
If $T$ is a $\beta$-SYT, then $S_3 (S_8 S_7 S_6 S_5 S_4 S_3) v_T = 0$.
\end{lemma}

\begin{proof}
$S_3$ commutes with $S_5, S_6, S_7, S_8$. So
\[
S_3 \cdot S_8 S_7 S_6 S_5 S_4 S_3 \;=\; S_8 S_7 S_6 S_5 \cdot S_3 S_4 S_3.
\]
We claim that for $T$ a $\beta$-SYT, the vector $S_3 S_4 S_3 v_T$ is supported on
the set $\{T, s_3 T\}$, and both have $5,6$ stacked in column $1$ (so $S_5$ kills).

\emph{Case (β1) $T(2,1)=3, T(1,3)=4$.} Axial distance $d_3 = 3$ between 3 and 4,
so $S_3 v_T = a_3 v_T + b_3 v_{s_3 T}$. In $s_3 T$, the entries 4 and 5 are now at
$(2,1)$ and $(3,1)$ — same column — so $S_4 v_{s_3 T} = 0$. Therefore
$S_4 S_3 v_T = a_3 S_4 v_T = a_3 (a_4 v_T + b_4 v_{s_4 T})$ where in $s_4 T$ the
entries 3 and 4 are at $(2,1)$ and $(3,1)$ — same column — so $S_3 v_{s_4 T} = 0$.
Therefore $S_3 S_4 S_3 v_T = a_3 a_4 (a_3 v_T + b_3 v_{s_3 T})$ —
supported on $\{T, s_3 T\}$.

\emph{Case (β1') $T(2,1)=3, T(2,2)=4$.} Here $d_3 = 1$ (3,4 in same row), so
$S_3 v_T = (q+1) v_T$. Then $S_3 S_4 S_3 v_T = (q+1) S_3 S_4 v_T$. Now
$d_4 = -2$, so $S_4 v_T = a_4 v_T + b_4 v_{s_4 T}$, and in $s_4 T$, 3 and 4 are
column-adjacent at $(2,1)$ and $(3,1)$ — so $S_3 v_{s_4 T} = 0$. Hence
$S_3 S_4 v_T = a_4 (q+1) v_T$, giving $S_3 S_4 S_3 v_T = -(q+1) v_T$ —
supported on $\{T\}$.

\emph{Case (β2) $T(1,3)=3, T(2,1)=4$.} Here 4 and 5 are column-adjacent at $(2,1),
(3,1)$ — so $S_4 v_T = 0$. Hence $S_4 S_3 v_T = a_3 \cdot 0 + b_3 S_4 v_{s_3 T}$.
In $s_3 T$, 3 is at $(2,1)$ and 4 is at $(1,3)$ (β1-like); $S_4 v_{s_3 T}$ produces
support $\{s_3 T, s_4 s_3 T\}$, and in $s_4 s_3 T$ the entries 3 and 4 are at $(2,1),
(3,1)$ — same column — so $S_3 v_{s_4 s_3 T} = 0$. Therefore $S_3 S_4 S_3 v_T =
b_3 a_4 (a_3' v_{s_3 T} + b_3' v_T)$ — supported on $\{T, s_3 T\}$.

In every case, the support is in $\{T, s_3 T\}$, both of which have $5$ at $(3,1)$
and $6$ at $(4,1)$ (the $\beta$-condition involves no entries $\le 4$, so $s_3$
preserves it). These are SC-at-5 SYTs: $S_5 v_T = S_5 v_{s_3 T} = 0$. Therefore
$S_8 S_7 S_6 S_5 \cdot (S_3 S_4 S_3) v_T = (S_8 S_7 S_6) \cdot 0 = 0$.
\end{proof}

\begin{lemma}[$L_\beta$] \label{lem:Lbeta}
If $T$ is a $\beta$-SYT, then $\Omega v_T = 0$.
\end{lemma}

\begin{proof}
We show $R'_8 R'_9 v_T \in \Vm$. By Lemma~\ref{lem:preserves}, $S_3, S_4, \ldots, S_8$ preserve
the $\Vp/\Vm$ split. Decompose $R'_9 v_T = w_+ + w_-$ with $w_\pm \in V_\pm$. Then
$R'_8 R'_9 v_T = R'_8 w_+ + R'_8 w_- = R'_8 w_+$ (since $R'_8(\Vm) = 0$).
So it suffices to show $\pi_+(R'_8 w_+) = 0$, where $\pi_+$ is projection to $\Vp$.

For SR-at-1 $T$, $S_1 v_T = (q+1) v_T$, so $R'_9 v_T = (q+1)(S_8 \cdots S_2) v_T$.
Tracking $\Vp/\Vm$: $S_2$ on SR-at-1 $T$ either stays in $\Vp$ (when 3 is in row 1
or column 2 of $T$) or splits into $\Vp \oplus \Vm$ where the $\Vm$ component
involves $s_2 T$ which is SC-at-1. After $S_2$, the operators $S_3, \ldots, S_8$
preserve the split. Therefore the $\Vp$-part of $R'_9 v_T$ is
$
w_+ \;=\; c_T \cdot (S_8 \cdots S_3) v_T
$
for some nonzero scalar $c_T$ depending on the position of 3 in $T$.

Now we claim $w_+ \in \Vpalpha = \Vp \cap \ker S_3$. By Lemma~\ref{lem:SLbeta},
$S_3 (S_8 \cdots S_3) v_T = 0$, so $S_3 w_+ = 0$. Combined with $w_+ \in \Vp$,
this gives $w_+ \in \Vpalpha$, i.e.\ $w_+ = \sum_{T' \in \alpha} c_{T'} v_{T'}$.

For each $\alpha$-SYT $T'$ in this sum, by $L_\alpha$ (Lemma~\ref{lem:Lalpha} with
$k=8$), $R'_8 v_{T'} \in \Vm$. Therefore $R'_8 w_+ \in \Vm$, completing the proof.
\end{proof}

\section{Proof of L\textsubscript{γ}}

We need an analogue of Lemma~\ref{lem:SLbeta} for $R'_8 v_{T'}$ (rather than
$R'_9$), and a finer support description.

\begin{lemma} \label{lem:SLbeta-R8}
If $T'$ is a $\beta$-SYT, then $S_3 (S_7 S_6 S_5 S_4 S_3) v_{T'} = 0$. Consequently
the $\Vp$-part of $R'_8 v_{T'}$ lies in $\Vpalpha$.
\end{lemma}

\begin{proof}
Verbatim analogue of Lemma~\ref{lem:SLbeta} with $S_8$ removed. The chain becomes
$S_3 \cdot S_7 S_6 S_5 S_4 S_3 = S_7 S_6 S_5 \cdot S_3 S_4 S_3$. The same case
analysis (β1, β1', β2) gives $S_3 S_4 S_3 v_{T'} \in \mathrm{span}\{v_{T'}, v_{s_3 T'}\}
\subseteq \ker S_5$. Hence the chain vanishes.

The second statement: as in Lemma~\ref{lem:Lbeta}, the $\Vp$-part of $R'_8 v_{T'}$
is a scalar times $(S_7 S_6 S_5 S_4 S_3) v_{T'}$, and the first statement gives
$S_3 \cdot$ this $= 0$, hence the $\Vp$-part is in $\Vpalpha$.
\end{proof}

\begin{lemma}[Support for γ] \label{lem:SLgamma}
If $T$ is a $\gamma$-SYT, then $\pi_+(R'_9 v_T) \in \Vpalpha \oplus \Vpbeta$.
\end{lemma}

\begin{remark} \emph{Status:}
This is the structural piece that the present proof does not yet fully close.
Computational verification at generic $q$ (file
\texttt{2026-05-19-gamma-structure.py}, output Q1) confirms the support claim
for all 6 $\gamma$-only SYTs of $\lambda = (3,3,1,1,1)$:
$\pi_+(R'_9 v_T)$ is supported on at most $8$ SYTs, all in $\alpha \cup \beta$.
The structural ingredients in place are:
\begin{itemize}
\item Each $s_3, s_4$ preserves the $\gamma$-condition (positions of $7,8$
unchanged), so $S_3 S_4 S_3 v_T$ is supported on $\gamma$-SYTs.
\item The chain $S_8 \cdots S_3$ acts within $\Vp$ (preserved by all $S_j$, $j \ge 3$).
\item For the subcases with $T(3,1) = 6$ (so $T$ has SC at 6 as well as 7), the
chain action restricts further; for $T(3,1) = 5$ the structure is more involved.
\end{itemize}
A full structural proof appears to require either a finer SC-at-7-propagation
argument through $S_5$ and $S_6$, or a direct algebraic identity in $H_n(q)$
expressing $S_3 \cdot (S_8 \cdots S_3) v_T$ in terms of $S_7 v_T$ for $\gamma$-SYTs.
\end{remark}

\begin{lemma}[$L_\gamma$] \label{lem:Lgamma}
If $T$ is a $\gamma$-SYT, then $\Omega v_T = 0$.
\end{lemma}

\begin{proof}[Proof, modulo Lemma~\ref{lem:SLgamma}]
We show $R'_7 R'_8 R'_9 v_T \in \Vm$.

Decompose $R'_9 v_T = w_+ + w_-$ with $w_- \in \Vm$. By Lemma~\ref{lem:SLgamma},
$w_+ = w_+^\alpha + w_+^\beta$ with $w_+^\alpha \in \Vpalpha$ and $w_+^\beta
\in \Vpbeta$.

Now $R'_8 R'_9 v_T = R'_8 w_+ + R'_8 w_-$. The $\Vm$ part is killed by $R'_8$:
$R'_8 w_- = 0$. So $R'_8 R'_9 v_T = R'_8 (w_+^\alpha + w_+^\beta)$.

For $w_+^\alpha = \sum c_{T'} v_{T'}$ over $\alpha$-SYTs $T'$: by Lemma~\ref{lem:Lalpha}
applied with $k=8$, each $R'_8 v_{T'} \in \Vm$. So $R'_8 w_+^\alpha \in \Vm$.

For $w_+^\beta = \sum d_{T'} v_{T'}$ over $\beta$-SYTs $T'$: by Lemma~\ref{lem:SLbeta-R8},
$\pi_+(R'_8 v_{T'}) \in \Vpalpha$. So $\pi_+(R'_8 w_+^\beta) \in \Vpalpha$.

Combining: $R'_8 R'_9 v_T \in \Vm \oplus \Vpalpha$.

Finally, $R'_7$ applied: $R'_7 \Vm = 0$ (rightmost $S_1$ kills). For the
$\Vpalpha$-part, decompose $\pi_+(R'_8 R'_9 v_T) = \sum e_{T'} v_{T'}$ over
$\alpha$-SYTs $T'$. By Lemma~\ref{lem:Lalpha} with $k=7$, $R'_7 v_{T'} \in \Vm$
for each. Hence $R'_7 (\pi_+(R'_8 R'_9 v_T)) \in \Vm$.

Therefore $R'_7 R'_8 R'_9 v_T \in \Vm$, and applying any further $R'_j$ gives 0.
\end{proof}

\section{Theorem proof}

\begin{proof}[Proof of Theorem~\ref{thm:main}]
($\Leftarrow$) If $T$ satisfies $\alpha$, $\beta$, or $\gamma$, then $\Omega v_T = 0$ by
Lemmas~\ref{lem:Lalpha}, \ref{lem:Lbeta}, \ref{lem:Lgamma} respectively.

($\Rightarrow$) Computational: verified at generic $q$ for all 50 SR-at-1 SYTs
(diagnostic in \texttt{2026-05-19-V33111-residual-diagnostic.py}). The
$15$ SR-at-1 SYTs not satisfying any of $\alpha, \beta, \gamma$ all have
$\Omega v_T \ne 0$.
\end{proof}

\section{Computational verification}

All three forward lemmas are verified at generic $q$ via SymPy
(\texttt{2026-05-19-key-lemma-verify.py}): $20/20$ $\alpha$-class, $9/9$
$\beta$-class, $6/6$ $\gamma$-class pass the test
$S_1 \cdot \Pi_{k+1}(T) = 0$ at the predicted first-vanish stage.

The support claims (the heart of the structural proofs) are also verified
(\texttt{2026-05-19-supports.py},
\texttt{2026-05-19-vplus-support.py},
\texttt{2026-05-19-gamma-structure.py}):
\begin{itemize}
\item For $\alpha$-SYT $T$, $R'_9 v_T \in \Vm$ (support entirely in SC-at-1).
\item For $\beta$-SYT $T$, $\pi_+(R'_9 v_T) \in \Vpalpha$
(support entirely in $\alpha$-SYTs).
\item For $\gamma$-SYT $T$, $\pi_+(R'_9 v_T) \in \Vpalpha \oplus \Vpbeta$
(support entirely in $\alpha \cup \beta$).
\end{itemize}

\section{Gap and outlook}

The remaining structural gap is Lemma~\ref{lem:SLgamma}: a $q$-generic proof that
for $\gamma$-SYT $T$, $\pi_+(R'_9 v_T)$ is supported in $\alpha \cup \beta$.
A natural attempt — the commutation reduction
$S_3 \cdot (S_8 \cdots S_3) = (S_8 S_7 S_6 S_5)(S_3 S_4 S_3)$ — reduces
to showing $(S_8 S_7 S_6 S_5)(S_3 S_4 S_3) v_T \in \Vm$ for $\gamma$-SYT,
but the $\gamma$-condition ($S_7 v_T = 0$) does not directly trigger here
because $s_5, s_6$ scramble entries 5,6,7 before the $S_7$ in the chain acts.

Two routes appear promising:
\begin{enumerate}
\item Track $\gamma$-class invariance through $S_5, S_6$: the support of
$(S_5 \cdot S_3 S_4 S_3) v_T$ stays within $\{T'  : T'(4,1)=7, T'(5,1)=8\}$,
but $(S_6 \cdot \text{that})$ may not. Need to compute how the $\Vp$-part
fragments under $S_6$ and check that the leftover $\Vp$-component
projects into $\Vpalpha \oplus \Vpbeta$ rather than the larger ``other'' subspace.
\item Build a direct Hecke-algebra identity of the form
$S_3 (S_8 \cdots S_3) = A \cdot S_7 + B$ where $B \in \Vpalpha \oplus \Vpbeta$
generation, applied to $v_T$ kills the first term via $S_7 v_T = 0$.
\end{enumerate}

The gap is finite — proving it for the six $\gamma$-only SYTs of
$\lambda=(3,3,1,1,1)$ closes the theorem for this shape — and the gap has been
verified computationally for all six.

\section*{Files referenced}
\begin{itemize}
\item Verification: \texttt{\textasciitilde/projects/scratch/2026-05-19-key-lemma-verify.py}
\item Support analysis: \texttt{\textasciitilde/projects/scratch/2026-05-19-supports.py},
\texttt{2026-05-19-vplus-support.py}, \texttt{2026-05-19-gamma-structure.py}
\item Empirical discriminator: \texttt{\textasciitilde/projects/scratch/2026-05-19-V33111-residual-diagnostic.md}
\end{itemize}

\end{document}
